% !TEX encoding = UTF-8
% !TEX Root = 0_main.tex

% \vspace{-0.2in}
\section{Preliminaries and Motivations}
% \vspace{-0.1in}

\noindent
\textbf{Generic adaptive methods.} Algorithm~\ref{algo:adaptive} is a generic framework (all operations are element-wise).
It describes various popular stochastic gradient descent algorithms~\citep{reddi2019convergence}.
% like Adam~\citep{kingma2014adam} and, Adagrad~\citep{duchi2011adaptive} and RMSprop~\citep{tieleman2012lecture}.
Specifically, different optimization algorithms can be specified by different choices of $\phi(.)$ and $\psi(.)$, where
$\phi(.)$ specifies how the momentum at time step $t$ is calculated, and $\psi(.)$ how the adaptive learning rate at $t$ is calculated.
For example, 
% in the vanilla SGD, these two functions are set to:
% \begin{align*}
% \vspace{-0.1in}
% \phi(g_1, \cdots, g_t) = g_t \quad\mbox{and}\quad \psi(g_1, \cdots, g_t) = 1.
% \vspace{-0.2in}
% \end{align*}
% And 
in the Adam algorithm, we have:
\begin{align}
% \vspace{-0.2in}
\phi(g_1, \cdots, g_t) = \frac{(1 - \beta_1)\sum_{i = 1}^t \beta_1^{t-i}g_i}{1 - \beta_1^t} 
\quad\mbox{and}\quad 
\psi(g_1, \cdots, g_t) = \sqrt{\frac{1-\beta_2^t}{(1 - \beta_2)\sum_{i = 1}^t \beta_2^{t-i}g_i^2}}.
\label{eqn:adam}
%\vspace{0.1in}
\end{align}
For numerical stability, the function $\psi(.)$ in Equation~\ref{eqn:adam} is usually calculated as $\hat{\psi}(g_1, \cdots, g_t) = \frac{\sqrt{1-\beta_2^t}}{\epsilon + \sqrt{(1 - \beta_2)\sum_{i = 1}^t \beta_2^{t-i}g_i^2}}$, 
where $\epsilon$ is a relatively small / negligible value (\eg, $1\times 10^{-8}$). 
% \vspace{-0.1in}
% \setlength{\textfloatsep}{-1pt}

\begin{algorithm}[!ht]
\DontPrintSemicolon
\KwIn{$\{\alpha_t\}_{t = 1}^T$: step size, $\{\phi_t, \psi_t\}_{t = 1}^T$: function to calculate momentum and adaptive rate, \newline
$\theta_0$: initial parameter, $f(\theta)$: stochastic objective function.}
\KwOut{$\theta_T$: resulting parameters}
\While{$t = 1$ to $T$}{
    $g_t \gets \nabla_{\theta} f_t(\theta_{t - 1})$ (Calculate gradients w.r.t. stochastic objective at timestep t)\;
    $m_t \gets \phi_t (g_1, \cdots, g_t)$ (Calculate momentum)\;
    $l_t \gets \psi_t (g_1, \cdots, g_t)$ (Calculate adaptive learning rate)\;
    $\theta_t \gets \theta_{t-1} - \alpha_t m_t l_t$ (Update parameters)\;
}
\Return{$\theta_T$}
\caption{Generic adaptive optimization method setup. All operations are element-wise. }
\label{algo:adaptive}
\end{algorithm}

% \begin{figure}[!htb]
\begin{figure}[t]
\centering
\vspace{-0.3in}
    \includegraphics[width=\textwidth]{fig/histogram_4.pdf}
\vspace{-0.6cm}
\caption{The absolute gradient histogram of the Transformers on the De-En IWSLT' 14 dataset during the training (stacked along the y-axis). X-axis is absolute value in the log scale and the height is the frequency. Without warmup, the gradient distribution is distorted in the first 10 steps. }
\vspace{-0.3cm}
\label{fig:histogram_2}
\end{figure}

\noindent
\textbf{Learning rate warmup.} Instead of setting the learning rate $\alpha_t$ as a constant or in a decreasing order, 
% a learning rate warmup strategy sets $\alpha_t$ as some small values in the first few steps.
a learning rate warmup strategy sets $\alpha_t$ as smaller values in the first few steps, thus not satisfying $\forall t\, \alpha_{t+1} \leq \alpha_{t}$.
% and decrease $\alpha_t$ after $T_w$ steps. 
For example, linear warmup sets $\alpha_t = t \,\alpha_0$ when $t < T_w$.
%In many machine learning applications, e.g., neural machine translation, it has been empirically observed that these adaptive stochastic optimization algorithms require a learning rate warmup stage, (\eg, a widely used warmup strategy is to set the learning rate as $\alpha_t = t \,\alpha_0$ in the first $T_w$ steps, \ie, $t < T_w$); otherwise, the loss is trapped in bad/suspicious local optima. S
Warmup has been demonstrated to be beneficial in many deep learning applications.
For example, in the NMT experiments in Figure~\ref{fig:eps2k}, the training loss convergences around 10 when warmup is not applied (Adam-vanilla), and it surprisingly decreases to below 3 after applying warmup (Adam-warmup). 


To further analyze this phenomenon, we visualize the histogram of the absolute value of gradients on a log scale in Figure~\ref{fig:histogram_2}.
We observe that, without applying warmup, the gradient distribution is distorted to have a mass center in relatively small values within 10 updates. Such gradient distortion means that the vanilla Adam is trapped in bad/suspicious local optima after the first few updates. 
Warmup essentially reduces the impact of these problematic updates to avoid the convergence problem.
%Intuitively, since the widely used adaptive learning rates are designed as the inverse scale of the gradient (\eg, $\psi(.)$ in Equation~\ref{eqn:adam}), small gradients would have larger adaptive learning rates, and random fluctuations of these gradients would be amplified and hinder the optimizer from advancing model training. 
% Therefore, the vanilla Adam will be trapped in this stage after the first few updates
%Therefore, such gradient distortion can trap the vanilla Adam after the first few updates, and it is necessary to fix those problematic updates to avoid the convergence problems.
In the following sections, we focus our analysis on learning rate warmup for the Adam algorithm, while it can be applied to other algorithms that use similar adaptive learning rate ($\psi(.)$) designs, \eg, RMSprop~\citep{tieleman2012lecture} and Nadam~\citep{dozat2016incorporating}. 


% \section{Preliminaries}

% As in Algorithm~\ref{algo:adaptive}, we follow the previous work~\citep{reddi2019convergence} and use a generic framework (all operations are element-wise) to describe adaptive gradient descent methods, which covers popular algorithms like Adam~\citep{kingma2014adam}, Adagrad~\citep{duchi2011adaptive} and RMSprop~\citep{tieleman2012lecture}.
% Thus, different optimization algorithms can be described as different choices of $\phi(.)$ and $\psi(.)$.
% For example, 
% % in the vanilla SGD, these two functions are set to:
% % $$
% % \phi(g_1, \cdots, g_t) = g_t \quad\mbox{and}\quad \psi(g_1, \cdots, g_t) = 1.
% % $$
% in the Adam algorithm, these two functions are set to:
% % $$
% % \phi(g_1, \cdots, g_t) = \frac{(1 - \beta_1)\sum_{i = 1}^t \beta_1^{t-i}g_t}{1 - \beta_1^t} 
% % \quad\mbox{and}\quad 
% % \psi(g_1, \cdots, g_t) = \frac{\sqrt{1-\beta_2^t}}{\sqrt{(1 - \beta_2)\sum_{i = 1}^t \beta_2^{t-i}g_i^2} + \epsilon}
% % $$
% % where $\epsilon$ is for numerical stability, and is usually set to relative small value (\eg, $1\times 10^{-8}$). In our most analysis, we 
% \begin{equation}
% \phi(g_1, \cdots, g_t) = \frac{(1 - \beta_1)\sum_{i = 1}^t \beta_1^{t-i}g_t}{1 - \beta_1^t} 
% \quad\mbox{and}\quad 
% \psi(g_1, \cdots, g_t) = \sqrt{\frac{1-\beta_2^t}{(1 - \beta_2)\sum_{i = 1}^t \beta_2^{t-i}g_i^2}}.
% \label{eqn:adam}
% \end{equation}
% For numerical stability, $\psi(.)$ is usually calculated as $\hat{\psi}(g_1, \cdots, g_t) = \frac{\sqrt{1-\beta_2^t}}{\epsilon + \sqrt{(1 - \beta_2)\sum_{i = 1}^t \beta_2^{t-i}g_i^2}}$, 
% where $\epsilon$ is set to relative small values (\eg, $1\times 10^{-8}$). % this introduction is used later to introduce Adam-eps.
% % We further summarize several popular optimization algorithms in Table~\ref{tab:adaptive}. 
% % It is worth mentioning that, all operations in this paper are element-wise. 

% \begin{algorithm}[t]
% \DontPrintSemicolon
% \KwIn{$\{\alpha_t\}_{t = 1}^T$: step size, 
% $\{\phi_t, \psi_t\}_{t = 1}^T$: function to calculate momentum and adaptive rate, \newline
% $\theta_0$: initial parameter, $f(\theta)$: stochastic objective function.}
% \KwOut{$\theta_t$: resulting parameters}
% \While{$t = 1$ to $T$}{
%     $g_t \gets \Delta_{\theta} f_t(\theta_{t - 1})$ (Calculate gradients w.r.t. stochastic objective at timestep t)\;
%     $m_t \gets \phi_t (g_1, \cdots, g_t)$ (Calculate momentum)\;
%     $v_t \gets \psi_t (g_1, \cdots, g_t)$ (Calculate adaptive learning rate)\;
%     $\theta_t \gets \theta_{t-1} - \alpha_t m_t v_t$ (Update parameters)\;
% }
% \Return{$\theta_T$}
% \caption{Generic Adaptive Method Setup. All operations are element-wise. }
% \label{algo:adaptive}
% \end{algorithm}

% In many machine learning applications, e.g., neural machine translation, it has been empirically observed that these adaptive stochastic optimization algorithms require a learning rate warmup stage, i.e., using small learning rates in the first few epochs; otherwise, the loss fails to descent normally or even explores. 
% For example, a widely used warmup strategy is to set learning rate as $\alpha_t = t \,\alpha_0$ in the first $T_w$ steps (\ie, $t < T_w$), 
% As in Figure~\ref{fig:eps2k}, the perplexity convergences around 500 when warmup is not employed (Adam-vanilla), and it successfully descents under 10 after applying warmup (Adam-warmup). 
% % We conduct further analysis on this phenomenon. 
% To further analyze this phenomenon, we visualize the histogram of the absolute value of gradients (on a log scale) in Figure~\ref{fig:histogram_2}.
% We observe that, without applying warm up, the gradient distribution is distorted to have a mass center in smaller values within 10 updates, and is trapped in this stage. 
% % Specifically, since small gradient would have larger adaptive learning rate, 
% % the variance of the gradient would be amplified and hinder the optimizer from advancing model training. 
% Intuitively, since the widely used adaptive learning rates are designed as the inverse scale of the gradient (\eg, $\psi(.)$ in Equation~\ref{eqn:adam}), small gradients would have larger adaptive learning rate, and their variance would be amplified and hinder the optimizer from advancing model training. 
% Therefore, we conjecture the convergence problems met by Adam-vanilla is caused by the problematic updates in the early stage.
% % that this issue is due to the lack of robustness of the adaptive learning rate. 
% % this issue is due to the lack of robustness of the adaptive learning rate. 
% % In the following section, we will first provide theoretical and empirical evidences to verify our intuition, then proceed to the design of our proposed rectification term. 
% In the following sections, we focus our analysis on the Adam algorithm, while it can be applied to other algorithms that uses similar adaptive learning rate ($\psi(.)$) designs, \eg, RMSprop~\citep{tieleman2012lecture} and Nadam~\citep{dozat2016incorporating}. 

% \begin{minipage}{0.35\textwidth}
% \begin{figure}[H]
% \centering
% % \begin{tabular}[H]{ c } 
% % \includegraphics[width=\textwidth]{fig/ppl.pdf}\\
% \includegraphics[width=\textwidth]{fig/ppl_loss_legend.pdf}
% % \end{tabular}
% \caption{Training of Transformers on the De-En IWSLT'14 dataset. Up: Training perplexity w.r.t. gradient update \#; Bottom: Training loss w.r.t. gradient update \#.}
% \label{fig:eps2k}
% \end{figure}
% \end{minipage}
% \;
% \begin{minipage}{0.64\textwidth}
% \begin{figure}[H]
% \centering
%     \includegraphics[width=\textwidth]{fig/histogram_2.pdf}
% \caption{The histogram of the absolute value of gradients (on a log scale) during the training of Transformers on the De-En IWSLT' 14 dataset. }
% \label{fig:histogram_2}
% \end{figure}
% \end{minipage}

% % \begin{figure}[t]
% % \centering
% %     \includegraphics[width=\textwidth]{fig/histogram_2.pdf}
% % \caption{The histogram of the absolute value of gradients (on a log scale) during the training of Transformers on the De-En IWSLT' 14 dataset. }
% % \label{fig:histogram_2}
% % \end{figure}

% % \begin{figure}[t]
% % \centering
% % \begin{subfigure}[b]{0.4\textwidth}
% %     \centering
% %     \includegraphics[width=\textwidth]{fig/ppl.pdf}
% %     % \caption{\,}
% %     % \label{subfig:eps2k_ppl}
% % \end{subfigure}%
% % \begin{subfigure}[b]{0.4\textwidth}
% %     \centering
% %     \includegraphics[width=\textwidth]{fig/loss.pdf}
% %     % \caption{\,}
% %     % \label{subfig:eps2k_loss}
% % \end{subfigure}
% % \caption{Training of Transformers on the De-En IWSLT'14 dataset. Left: Training perplexity in the first 50k gradient updates; Right: Training loss in the first 50k gradient updates.}
% %     \label{fig:eps2k}
% % \end{figure}







% % \begin{figure}[t]
% % \centering
% %     \includegraphics[width=\textwidth]{fig/histogram_2.pdf}
% % \caption{The histogram of the absolute value of gradients during the training of Transformers on the De-En IWSLT' 14 dataset. }
% % \label{fig:histogram_2}
% % \end{figure}

% % \begin{table}[t]
% % \begin{center}
% % \caption{The choices of $\phi(.)$ and $\psi(.)$ of popular optimization algorithms.}
% % \label{tab:adaptive}
% % \begin{tabularx}{\columnwidth}{r *{2}{Y}}
% % \toprule
% %         & $\phi(g_1, \cdots, g_t)$ & $\psi(g_1, \cdots, g_t)$ \\\midrule
% % Vanilla SGD & $g_t$ & $1$ \\ \midrule

% % Adagrad & $g_t$ & $\sqrt{\frac{1}{\sum_{i = 1}^t g_i^2}}$\\\midrule

% % Adam & $\frac{(1 - \beta_1)\sum_{i = 1}^t \beta_1^{t-i}g_t}{1 - \beta_1^t}$ & $\sqrt{\frac{1-\beta_2^t}{(1 - \beta_2)\sum_{i = 1}^t \beta_2^{t-i}g_i^2}}$\\\midrule

% % RMSprop & $g_t$ & $\sqrt{\frac{1}{(1 - \beta_2)\sum_{i = 1}^t \beta_2^{t-i}g_i^2}}$\\
% % \bottomrule
% % \end{tabularx}
% % \end{center}
% % \end{table}
